{{دیوان
| عنوان = دُعا (ہے یہی میری نماز، ہے یہی میرا وضو)
| ماخذ = ویکی ماخذ
| ماخذ_لنک = https://ur.wikisource.org/wiki/%D8%AF%D9%8F%D8%B9%D8%A7_%28%DB%81%DB%92_%DB%8C%DB%81%DB%8C_%D9%85%DB%8C%D8%B1%DB%8C_%D9%86%D9%85%D8%A7%D8%B2%D8%8C_%DB%81%DB%92_%DB%8C%DB%81%DB%8C_%D9%85%DB%8C%D8%B1%D8%A7_%D9%88%D8%B6%D9%88%29
| مقابلہ = اقبال سائبر لائبریری (اقبال اکادمی پاکستان)
| مقابلہ_لنک = https://www.iqbalcyberlibrary.net
}}
=== (مسجدِ قُرطُبہ میں لِکھّی گئی) ===

<poem>
ہے یہی میری نماز، ہے یہی میرا وضو
میری نواؤں میں ہے میرے جگر کا لہُو

صُحبتِ اہلِ صفا، نُور و حضور و سُرور
سر خوش و پُرسوز ہے لالہ لبِ آبُجو

راہِ محبّت میں ہے کون کسی کا رفیق
ساتھ مرے رہ گئی ایک مری آرزو

میرا نشیمن نہیں درگہِ میر و وزیر
میرا نشیمن بھی تُو، شاخِ نشیمن بھی تُو

تجھ سے گریباں مرا مطلعِ صُبحِ نشور
تجھ سے مرے سینے میں آتشِ ’اَﷲ ھو‘

تجھ سے مری زندگی سوز و تب و درد و داغ
تُو ہی مری آرزو، تُو ہی مری جُستجو

پاس اگر تُو نہیں، شہر ہے ویراں تمام
تُو ہے تو آباد ہیں اُجڑے ہُوئے کاخ و کُو

پھر وہ شرابِ کُہن مجھ کو عطا کر کہ مَیں
ڈھُونڈ رہا ہوں اُسے توڑ کے جام و سُبو

چشمِ کرم ساقیا! دیر سے ہیں منتظر
جلوَتیوں کے سُبو، خلوَتیوں کے کُدو

تیری خدائی سے ہے میرے جُنوں کو گِلہ
اپنے لیے لامکاں، میرے لیے چار سُو!

فلسفہ و شعر کی اور حقیقت ہے کیا
حرفِ تمنّا، جسے کہہ نہ سکیں رُو بُرو
</poem>
